% ═══════════════════════════════════════════════════════════════
% AB-04: Stromdetektive im Haushalt
% Physik Klasse 8 | Projektarbeit
% ═══════════════════════════════════════════════════════════════

\documentclass[11pt, a4paper]{article}

\usepackage[utf8]{inputenc}
\usepackage[T1]{fontenc}
\usepackage[ngerman]{babel}

\usepackage{helvet}
\renewcommand{\familydefault}{\sfdefault}

\usepackage{microtype}
\usepackage[margin=1.5cm, top=1.8cm, bottom=1.5cm]{geometry}
\usepackage{enumitem}
\usepackage{tabularx}
\usepackage{booktabs}
\usepackage{array}
\usepackage{multirow}
\usepackage{colortbl}
\usepackage{tikz}
\usetikzlibrary{calc}

\usepackage{amsmath, amssymb}
\usepackage{siunitx}
\sisetup{locale = DE, per-mode = symbol}

\usepackage{fancyhdr}
\usepackage{lastpage}
\usepackage{needspace}
\usepackage{graphicx}
\usepackage[hidelinks]{hyperref}
\usepackage{qrcode}

\usepackage{xcolor}
\usepackage{tcolorbox}
\tcbuselibrary{skins, breakable}

% Farben
\definecolor{physik}{HTML}{2c5aa0}
\definecolor{hellgrau}{HTML}{F5F5F5}
\definecolor{mittelgrau}{HTML}{888888}
\definecolor{dunkelgrau}{HTML}{444444}
\definecolor{rahmen}{HTML}{CCCCCC}
\definecolor{energie}{HTML}{FF9800}
\definecolor{geld}{HTML}{4CAF50}

% Variablen
\newcommand{\fach}{Physik}
\newcommand{\thema}{Stromdetektive im Haushalt}
\newcommand{\klasse}{8}
\newcommand{\maxpunkte}{25}

\colorlet{fachfarbe}{physik}

% Kopf- und Fußzeile
\pagestyle{fancy}
\fancyhf{}
\renewcommand{\headrulewidth}{0.4pt}
\lhead{\small \fach{} Klasse \klasse{} | \thema}
\rhead{\small Namen: \underline{\hspace{5cm}}}
\cfoot{\small\textcolor{mittelgrau}{Seite \thepage{} von \pageref{LastPage}}}

% Boxen
\newtcolorbox{merkkasten}[1][]{
    colback=hellgrau,
    colframe=hellgrau,
    boxrule=0pt,
    arc=0pt,
    left=8pt, right=8pt, top=8pt, bottom=6pt,
    borderline north={2pt}{0pt}{fachfarbe},
    before skip=8pt, after skip=8pt,
    #1
}

\newtcolorbox{formelkasten}{
    colback=white,
    colframe=fachfarbe,
    boxrule=1.5pt,
    arc=0pt,
    left=8pt, right=8pt, top=6pt, bottom=6pt,
    before skip=6pt, after skip=6pt
}

\newtcolorbox{hinweiskasten}{
    colback=white,
    colframe=white,
    boxrule=0pt,
    arc=0pt,
    left=8pt, right=4pt, top=4pt, bottom=4pt,
    borderline west={3pt}{0pt}{fachfarbe}
}

\newtcolorbox{tippkasten}{
    colback=white,
    colframe=energie,
    boxrule=1.5pt,
    arc=4pt,
    left=8pt, right=8pt, top=6pt, bottom=6pt,
    before skip=6pt, after skip=6pt
}

% Befehle
\newcommand{\luecke}[1]{\underline{\hspace{#1}}}
\newcommand{\punktebox}[1]{\hfill\small\textcolor{mittelgrau}{[\hspace{0.8em}/#1]}}

\newcommand{\aufgabe}[2]{%
    \needspace{4\baselineskip}%
    \vspace{10pt}%
    \noindent\textbf{\textcolor{fachfarbe}{Aufgabe #1}} #2%
    \vspace{4pt}%
    \nopagebreak[4]%
}

\newcommand{\operator}[1]{\textbf{#1}}

\newlist{teil}{enumerate}{2}
\setlist[teil,1]{label=\alph*), leftmargin=1.8em, itemsep=4pt, parsep=0pt, topsep=4pt}

\newcommand{\antwortzeilen}[1]{%
    \par\vspace{4pt}%
    \foreach \n in {1,...,#1}{%
        \noindent\rule{\linewidth}{0.3pt}\vspace{8pt}%
    }%
}

\setlength{\parindent}{0pt}
\setlength{\parskip}{3pt}

\begin{document}

% ═══════════════════════════════════════════════════════════════
% HEADER
% ═══════════════════════════════════════════════════════════════

\noindent
\begin{minipage}[c]{0.85\textwidth}
    \begin{center}
        {\Large\bfseries\textcolor{fachfarbe}{\thema}}\\[0.2em]
        {\normalsize Partnerarbeit}
    \end{center}
\end{minipage}%
\hfill%
\begin{minipage}[c]{0.12\textwidth}
    \raggedleft
    \href{https://devbox42.github.io/sciencesim/physik/kl08/04-stromdetektive/LP-04-stromdetektive.html}{%
        \qrcode[height=1.5cm]{https://devbox42.github.io/sciencesim/physik/kl08/04-stromdetektive/LP-04-stromdetektive.html}}
\end{minipage}

\vspace{-0.2em}
\hrule
\vspace{0.5em}

% ═══════════════════════════════════════════════════════════════
% INFOBOX
% ═══════════════════════════════════════════════════════════════

\begin{merkkasten}
\small
\textbf{Aufgabe:} Ihr seid Stromdetektive! Analysiert den Stromverbrauch in eurem Haushalt.
\hfill
\textbf{Formeln:} \fbox{$E = P \cdot t$} \quad \fbox{Kosten $= E \cdot$ Strompreis}

\vspace{0.2em}
\begin{tabular}{@{}ll@{\hspace{1.5cm}}ll@{\hspace{1.5cm}}l@{}}
$E$ = Energie (Wh/kWh) & $P$ = Leistung (W/kW) & $t$ = Zeit (h) & Strompreis: \textbf{30 ct/kWh}
\end{tabular}
\hfill \textbf{Umrechnung:} 1\,kWh = 1000\,Wh
\end{merkkasten}

% ═══════════════════════════════════════════════════════════════
% TABELLE (Zwei-Spalten-Layout, 10 Zeilen pro Raum)
% ═══════════════════════════════════════════════════════════════

\vspace{0.1em}
\textbf{\textcolor{physik}{Geräte-Erfassung:}} Geht Raum für Raum durch euren Haushalt. \punktebox{5}

\vspace{0.2em}
\scriptsize
\renewcommand{\arraystretch}{0.82}
\setlength{\tabcolsep}{2pt}

\newcommand{\roomtable}[1]{%
\begin{tabular}{|p{1.75cm}|p{0.75cm}|p{0.75cm}|p{0.9cm}|p{0.75cm}|p{0.9cm}|p{0.9cm}|}
\hline
\multicolumn{7}{|c|}{\cellcolor{physik!20}\textbf{#1}} \\
\hline
\textbf{Gerät} & \textbf{P} \newline \textit{W} & \textbf{t} \newline \textit{h/Tag} & \textbf{E} \newline \textit{Wh/Tag} & \textbf{Tage} \newline \textit{/Jahr} & \textbf{E} \newline \textit{kWh/J} & \textbf{Kosten} \newline \textit{€/Jahr} \\
\hline
 & & & & & & \\ \hline
 & & & & & & \\ \hline
 & & & & & & \\ \hline
 & & & & & & \\ \hline
 & & & & & & \\ \hline
 & & & & & & \\ \hline
 & & & & & & \\ \hline
 & & & & & & \\ \hline
 & & & & & & \\ \hline
 & & & & & & \\ \hline
\end{tabular}%
}

\noindent
\begin{minipage}[t]{0.48\textwidth}
\centering
\roomtable{Küche}
\vspace{0.15em}

\roomtable{Wohnzimmer}
\vspace{0.15em}

\roomtable{Schlafzimmer}
\vspace{0.15em}

\roomtable{Keller / Garage}
\end{minipage}%
\hfill%
\begin{minipage}[t]{0.48\textwidth}
\centering
\roomtable{Bad}
\vspace{0.15em}

\roomtable{Arbeitszimmer}
\vspace{0.15em}

\roomtable{Kinderzimmer}
\vspace{0.15em}

\roomtable{Flur / Sonstige}
\end{minipage}

\normalsize
\setlength{\tabcolsep}{6pt}
\vspace{0.15em}
\begin{tippkasten}
\scriptsize
\textbf{Formeln:} $E_{\text{Tag}} = P \cdot t$ \quad $E_{\text{Jahr}} = E_{\text{Tag}} \cdot \text{Tage} \div 1000$ \quad $\text{Kosten} = E_{\text{Jahr}} \cdot 0{,}30\,\text{€}$ \quad\quad \textbf{Werte:} LED 10\,W | TV 100\,W | Föhn 2000\,W
\end{tippkasten}

% ═══════════════════════════════════════════════════════════════
% SEITE 2
% ═══════════════════════════════════════════════════════════════

\newpage

% ═══════════════════════════════════════════════════════════════
% AUFGABE 1: TOP 3
% ═══════════════════════════════════════════════════════════════

\aufgabe{1}{– Die größten Verbraucher} \punktebox{4}

\operator{Identifiziere} die drei Geräte mit dem höchsten täglichen Energieverbrauch aus deiner Tabelle.

\vspace{0.3em}
\small
\renewcommand{\arraystretch}{1.2}
\begin{tabular}{|c|p{4.5cm}|c|p{5cm}|}
\hline
\textbf{Platz} & \textbf{Gerät} & \textbf{$E$ (Wh/Tag)} & \textbf{Warum so viel?} \\
\hline
1 & & & \\
\hline
2 & & & \\
\hline
3 & & & \\
\hline
\end{tabular}
\normalsize

% ═══════════════════════════════════════════════════════════════
% AUFGABE 2: KOSTEN BERECHNEN
% ═══════════════════════════════════════════════════════════════

\aufgabe{2}{– Stromkosten berechnen} \punktebox{7}

\begin{teil}
    \item \operator{Berechne} den Tagesverbrauch (Summe aller Geräte). \hfill (2 BE)
    \quad $E_{\text{Tag}} = $ \luecke{2.5cm} Wh = \luecke{1.8cm} kWh

    \item \operator{Berechne} die monatlichen Stromkosten (30 Tage, 30 ct/kWh). \hfill (3 BE)

    \antwortzeilen{1}
    Kosten pro Monat: \luecke{2.5cm} €

    \item \operator{Berechne} die Jahreskosten (365 Tage). \hfill (2 BE)
    \quad Kosten pro Jahr: \luecke{2.5cm} €
\end{teil}

% ═══════════════════════════════════════════════════════════════
% AUFGABE 3: SPARMASSNAHMEN
% ═══════════════════════════════════════════════════════════════

\aufgabe{3}{– Sparmaßnahmen} \punktebox{5}

\operator{Nenne} für eure drei größten Verbraucher je eine Sparmaßnahme. \textit{\small (z.\,B. kürzere Nutzung, niedrigere Temperatur, LED statt Glühbirne, Standby aus)}

\vspace{0.3em}
\small
\renewcommand{\arraystretch}{1.2}
\begin{tabular}{|p{3.5cm}|p{11cm}|}
\hline
\textbf{Gerät} & \textbf{Sparmaßnahme} \\
\hline
& \\
\hline
& \\
\hline
& \\
\hline
\end{tabular}
\normalsize

% ═══════════════════════════════════════════════════════════════
% AUFGABE 4: EINSPARPOTENZIAL
% ═══════════════════════════════════════════════════════════════

\aufgabe{4}{– Einsparpotenzial berechnen} \punktebox{4}

\operator{Wähle} eine Sparmaßnahme und \operator{berechne}, wie viel ihr pro Jahr sparen könntet.

\vspace{0.2em}
Gewählte Maßnahme: \luecke{10cm}

\vspace{0.2em}
\begin{tabular}{@{}p{7.5cm}p{7.5cm}@{}}
\textbf{Vorher:} & \textbf{Nachher:} \\
$P$: \luecke{1.5cm} W \quad $t$: \luecke{1.5cm} h/Tag & $P$: \luecke{1.5cm} W \quad $t$: \luecke{1.5cm} h/Tag \\
$E$: \luecke{2cm} Wh/Tag & $E$: \luecke{2cm} Wh/Tag
\end{tabular}

\vspace{0.2em}
\textbf{Ersparnis:} $\Delta E_{\text{Tag}} = $ \luecke{2cm} Wh/Tag \quad $\Delta E_{\text{Jahr}} = \Delta E \cdot 365 = $ \luecke{2cm} kWh/Jahr

\vspace{0.2em}
\textbf{Geldersparnis:} $\Delta E_{\text{Jahr}} \cdot 0{,}30\,\text{€/kWh} = $ \luecke{2.5cm} \textbf{€/Jahr}

\vspace{0.3em}
\begin{formelkasten}
\centering
\textbf{Unser Spar-Tipp:} Mit \textit{``\luecke{4cm}''} sparen wir ca. \luecke{1.8cm} €/Jahr!
\end{formelkasten}

\end{document}
